% Background

\section{Background}
\label{sec:background}

\subsection{eBPF for Security Monitoring}
\label{sec:bg-ebpf}

eBPF (extended Berkeley Packet Filter) enables user-defined programs to run safely within the Linux kernel~\cite{ebpf-intro}.
Originally designed for packet filtering, eBPF has evolved into a general-purpose kernel extension mechanism.
Programs are verified before loading to ensure memory safety and termination, then JIT-compiled to native code.

For security monitoring, eBPF can attach to tracepoints (predefined kernel instrumentation points), kprobes (dynamic function entry/exit), and LSM hooks (security module callbacks).
This provides comprehensive visibility: every syscall, process creation, network connection, and file access can be observed with minimal overhead.

Captured events must be transferred to userspace for analysis.
eBPF offers two mechanisms: \textit{perf buffers} provide per-CPU queues with copy semantics, while \textit{ring buffers} (Linux 5.8+) enable memory-mapped access with better cache behavior.
\textsc{WasmSentinel} uses ring buffers to enable our zero-copy data channel.

\subsection{WebAssembly Sandboxing}
\label{sec:bg-wasm}

WebAssembly (Wasm) is a portable binary format originally designed for web browsers but increasingly used for server-side sandboxing~\cite{wasm-spec}.
Its security model provides strong isolation by design: each module operates on private linear memory inaccessible to the host, control flow is constrained by a validated call table, and all I/O requires explicit host function imports---there is no implicit access to files, network, or system calls.

Modern Wasm runtimes like Wasmtime~\cite{wasmtime} achieve near-native performance (5--20\% overhead) through ahead-of-time compilation.
WASI (WebAssembly System Interface)~\cite{wasi} standardizes capability-based access to OS resources, allowing fine-grained control over what each module can access.

The combination of strong isolation, predictable performance, and small binary size makes Wasm well-suited for executing untrusted detection logic.

\subsection{Syscall-Based Threat Detection}
\label{sec:bg-ml}

System call sequences have long been recognized as effective indicators of malicious behavior~\cite{stide}.
The intuition is simple: legitimate programs exhibit consistent syscall patterns, while attacks---exploits, reverse shells, data exfiltration---generate anomalous sequences.

STIDE (Sequence Time-Delay Embedding) pioneered n-gram analysis: build a database of ``normal'' syscall n-grams during training, then flag sequences containing novel n-grams as anomalous.
This approach is computationally efficient but has limited precision, as unusual but benign workloads can trigger false positives.
Neural network classifiers can learn more nuanced patterns but typically require heavyweight frameworks (TensorFlow, PyTorch) with millisecond-scale latency and no isolation guarantees.

\textsc{WasmSentinel} bridges this gap: we implement lightweight ML models directly in Rust, compile to Wasm, and achieve microsecond inference within a sandbox.
